LiionTR is currently a reduced-order thermal runaway framework with
explicit models for reaction kinetics, heat generation, gas generation,
pressure, and venting.

Several extensions are required before the framework can represent the
full range of physical processes relevant to lithium-ion battery failure.

Development priorities are described below.


\subsection{Thermochemical gas modeling}

The current gas-generation formulation uses prescribed species yields.

This approach is useful for empirical models and verification but cannot
predict gas composition from the chemical state of the reacting battery
materials.

A future thermochemical formulation will represent reaction products in
terms of elemental inventories.

Conceptually,

\begin{equation}
    \text{thermal decomposition}
    \rightarrow
    \text{element source terms}
    \rightarrow
    \text{element inventory}
    \rightarrow
    \text{thermochemical equilibrium}.
\end{equation}

Relevant elements may include

\begin{equation}
    \mathrm{
        C,\,
        H,\,
        O,\,
        N,\,
        F,\,
        P,\,
        Li
    }.
\end{equation}

A thermochemical equilibrium backend can then determine gas composition
subject to specified thermodynamic constraints.


\subsection{Cantera integration}

Cantera is a strong candidate for the thermochemical backend.

A future LiionTR interface may provide a generic abstraction such as

\begin{equation}
    \texttt{GasEquilibriumBackend},
\end{equation}

with a Cantera implementation responsible for

\begin{itemize}
    \item equilibrium gas composition;
    \item mixture molar mass;
    \item heat capacities;
    \item enthalpy and internal energy;
    \item heat-capacity ratio;
    \item other thermodynamic properties required by venting.
\end{itemize}

The initial implementation may use constant-volume, constant-temperature
equilibrium calculations.

More advanced formulations may require coupled energy-conserving
equilibrium calculations.


\subsection{Energy-consistent thermochemistry}

The current thermal state is temperature.

This is convenient for reduced-order calculations but may become
insufficient when significant energy is redistributed through chemical
equilibrium, vaporization, gas discharge, or phase change.

A future formulation may therefore evolve internal energy,

\begin{equation}
    U,
\end{equation}

rather than temperature directly.

Temperature would then be reconstructed from the thermodynamic state,

\begin{equation}
    T
    =
    T
    \left(
        U,
        V,
        \mathbf{n}
    \right).
\end{equation}

Such a formulation would provide a more rigorous basis for energetic
coupling between condensed reactions and gas-phase thermochemistry.


\subsection{Electrochemical coupling}

Thermal runaway behavior depends strongly on the electrochemical state
of the cell.

Relevant quantities include

\begin{itemize}
    \item state of charge;
    \item electrode lithium concentration;
    \item electrode potential;
    \item current;
    \item irreversible electrochemical heating;
    \item reversible entropic heating.
\end{itemize}

PyBaMM may eventually provide these quantities to LiionTR.

The intended coupling is not to replace PyBaMM electrochemistry, but to
allow electrochemical state variables to modify thermal runaway kinetics,
reaction inventories, or onset conditions.


\subsection{Calorimetric parameter estimation}

Thermal runaway kinetic parameters are often derived from differential
scanning calorimetry or accelerating rate calorimetry measurements.

LiionTR may therefore include tools for fitting model parameters to
experimental data.

Potential parameters include

\begin{equation}
    A,
    \qquad
    E_a,
    \qquad
    \Delta H,
    \qquad
    n,
    \qquad
    \alpha_0.
\end{equation}

Parameter estimation could be formulated as an optimization problem,

\begin{equation}
    \boldsymbol{\theta}^{*}
    =
    \arg\min_{\boldsymbol{\theta}}
    J
    \left(
        \boldsymbol{\theta}
    \right),
\end{equation}

where $\boldsymbol{\theta}$ is the parameter vector and $J$ measures the
difference between simulation and experiment.


\subsection{Sensitivity and uncertainty analysis}

Published thermal runaway parameters can vary substantially between
cells, chemistries, measurement methods, and states of charge.

Future development should therefore consider parameter uncertainty.

Relevant outputs may include sensitivities of

\begin{itemize}
    \item thermal runaway onset time;
    \item maximum temperature;
    \item maximum pressure;
    \item vent-opening time;
    \item total gas generation;
    \item total energy release.
\end{itemize}

Both local sensitivity analysis and Monte Carlo uncertainty propagation
may be useful.


\subsection{Temperature-dependent material properties}

The current cell thermal properties are effectively constant.

Future material models may support

\begin{equation}
    c_p
    =
    c_p(T),
\end{equation}

\begin{equation}
    k
    =
    k(T),
\end{equation}

and potentially

\begin{equation}
    \rho
    =
    \rho(T).
\end{equation}

Phase transitions and latent-heat effects may also be relevant at severe
thermal runaway temperatures.


\subsection{Spatial thermal models}

The present lumped formulation cannot reproduce internal temperature
gradients.

A natural extension is a one-dimensional radial model for cylindrical
cells,

\begin{equation}
    \rho c_p
    \frac{\partial T}{\partial t}
    =
    \frac{1}{r}
    \frac{\partial}{\partial r}
    \left(
        k r
        \frac{\partial T}{\partial r}
    \right)
    +
    \dot{q}_{\mathrm{gen}}.
\end{equation}

More general finite-volume or finite-element formulations could later
support cylindrical, prismatic, and pouch cells.


\subsection{Cell-to-cell propagation}

Battery-module safety analysis requires coupling between multiple cells.

For cell $i$, a generic thermal balance may be written as

\begin{equation}
    C_i
    \frac{dT_i}{dt}
    =
    \dot{Q}_{i,\mathrm{gen}}
    -
    \dot{Q}_{i,\mathrm{ambient}}
    +
    \sum_j
    \dot{Q}_{j\rightarrow i}.
\end{equation}

The inter-cell term

\begin{equation}
    \dot{Q}_{j\rightarrow i}
\end{equation}

may include conduction, radiation, convection, or heat transfer from
vent gases and flames.

The modular cell representation already provides a basis for this
extension.


\subsection{Mechanical venting and rupture}

The current vent-opening model uses a prescribed pressure threshold and
instantaneous transition to a fixed vent area.

More realistic models may include

\begin{itemize}
    \item progressive vent displacement;
    \item pressure-dependent vent area;
    \item burst-disk failure;
    \item casing deformation;
    \item rupture;
    \item multiple vent paths.
\end{itemize}

A mechanical model could therefore determine

\begin{equation}
    A_v
    =
    A_v
    \left(
        P,
        T,
        t
    \right)
\end{equation}

instead of treating vent area as constant.


\subsection{Multiphase venting}

Real battery vent products may contain gases, vapor, droplets, aerosols,
and solid particles.

The current single-phase ideal-gas discharge model does not represent
these phenomena.

Future formulations may require

\begin{itemize}
    \item electrolyte vaporization;
    \item flashing flow;
    \item liquid entrainment;
    \item aerosol transport;
    \item multiphase choking criteria.
\end{itemize}


\subsection{External vent-jet and combustion models}

Once material leaves the cell, additional physical processes may occur
outside the battery.

These include

\begin{itemize}
    \item turbulent jet formation;
    \item mixing with ambient air;
    \item ignition;
    \item flame propagation;
    \item radiant heat transfer;
    \item toxic-species transport.
\end{itemize}

These processes are outside the current LiionTR domain but could
eventually be represented through interfaces to combustion or
computational-fluid-dynamics tools.


\subsection{Validation database}

A robust thermal runaway framework requires systematic comparison with
experimental measurements.

A future validation database could contain

\begin{itemize}
    \item battery-cell metadata;
    \item chemistry;
    \item geometry;
    \item nominal capacity;
    \item state of charge;
    \item abuse condition;
    \item ARC data;
    \item DSC data;
    \item temperature histories;
    \item pressure histories;
    \item vent gas composition;
    \item vented mass;
    \item photographs or high-speed imaging metadata.
\end{itemize}

Standardized datasets would permit repeatable comparison between
alternative LiionTR models.


\subsection{Software development}

In addition to physical-model development, several software improvements
are anticipated:

\begin{itemize}
    \item stable public APIs;
    \item structured serialization of simulations;
    \item model configuration files;
    \item reproducible simulation metadata;
    \item plotting and post-processing utilities;
    \item parameter databases;
    \item expanded example simulations;
    \item automated benchmarking;
    \item optional parallel parameter studies.
\end{itemize}


\subsection{Long-term architecture}

The long-term architecture can be summarized conceptually as

\begin{equation}
    \boxed{
    \begin{array}{c}
        \text{Electrochemistry} \\
        \text{PyBaMM}
    \end{array}
    }
    \longrightarrow
    \boxed{
    \begin{array}{c}
        \text{Thermal runaway} \\
        \text{LiionTR}
    \end{array}
    }
    \longrightarrow
    \boxed{
    \begin{array}{c}
        \text{Thermochemistry} \\
        \text{Cantera}
    \end{array}
    }.
\end{equation}

LiionTR is intended to remain the layer responsible for coordinating
battery-safety-specific physics, thermal runaway reactions, gas
generation, pressure, venting, and propagation.

The framework should therefore remain modular enough that increasing the
fidelity of one physical subsystem does not require redesigning the
complete simulation architecture.